% !TeX root = ../../tesis.tex

\subsection{Modelo de Datos y Esquema Relacional}
\label{subsec:modelo-datos}

\subsubsection{Esquema Entidad-Relación E-R}
\label{subsubsec:EsquemaER}

El diseño del esquema relacional busca simplificar el estándar complejo del
\gtfs{} (que incluye tablas como \textit{trips}, \textit{stop\_times} y
\textit{shapes}) en tres entidades primordiales optimizadas para el consumo de
la \api{} y el cálculo de indicadores topológicos:

\begin{figure}[H]
    \centering
    \includegraphics[width=0.85\textwidth]{Figures/Cap4/DiagramaER.png}
    \caption[Esquema Relacional Simplificado de Apimetro]{Representación de las entidades principales (Estaciones, Líneas y Polígonos) y sus relaciones para el cálculo de indicadores topológicos.}
    \label{fig:diagrama-er}
    \fuentefigura{Elaboración propia.}
\end{figure}


Como se observa en la Figura \ref{fig:diagrama-er}, el modelo se articula bajo
los siguientes componentes:

\begin{itemize}
\item \textbf{Estaciones:} Nodo puntual que contiene atributos de conectividad y
ubicación.
\item \textbf{Líneas:} Aristas de tipo \texttt{MultiLineString} que almacenan
métricas operativas (velocidad, capacidad, frecuencia).
\item \textbf{Polígonos:} Áreas administrativas para filtrado por alcaldía o
municipio.
\item \textbf{Relaciones:} Las líneas se vinculan con las estaciones a través de
un identificador de sistema y ramal, permitiendo la reconstrucción del grafo
metropolitano.
\end{itemize}


\subsubsection{Definición de Tablas Principales}
\label{subsubsec:DefinicionSQL}

El esquema relacional se implementa en tres tablas: \texttt{estaciones} (nodos
puntuales con atributos de conectividad), \texttt{lineas} (aristas
\texttt{MultiLineString} con métricas operativas de velocidad, frecuencia y
capacidad) y \texttt{poligonos} (áreas administrativas para filtrado por
alcaldía o municipio). Todas las columnas de geometría utilizan el tipo
\texttt{GEOMETRY} con proyección WGS~84 (SRID~4326) e índices espaciales
\texttt{GIST} para optimizar las consultas de ruteo. La definición completa del
esquema \textsc{sql} se presenta en el Anexo~\ref{anx:apimetro-sql}.

La arquitectura de estas tablas responde a dos necesidades fundamentales
identificadas en la investigación:
\begin{itemize}
\item \textbf{Optimización de Consultas Geográficas:} El uso de columnas de tipo
\texttt{GEOMETRY} y la creación de índices
\texttt{GIST} permiten realizar operaciones espaciales masivas en milisegundos,
tales como \texttt{ST\_AsGeoJSON} para la
exportación de geometrías, \texttt{ST\_DWithin} para el cálculo de áreas de
influencia o cobertura, y \texttt{ST\_Length} para
    determinar las distancias reales sobre la red vial.
\item \textbf{Cálculo de Indicadores Operativos:} A diferencia de un
\texttt{SIG}
convencional, la tabla \texttt{lineas} integra atributos operativos críticos
como la \textit{capacidad} del vehículo y la
\textit{frecuencia} de paso. Esto permite que el motor \texttt{VFTModel} calcule
automáticamente indicadores de desempeño como la
\textbf{Fuerza Nodal ($C_i$)} y el \textbf{Tiempo de Viaje ($T$)}, ponderando el
grafo metropolitano no solo por su fricción espacial,
sino por la oferta real de transporte y las penalizaciones por transferencia.
\end{itemize}

Finalmente, es importante señalar que esta infraestructura está alineada con los
preceptos de soberanía tecnológica y reproducibilidad científica, siendo
totalmente instalable en cualquier entorno (Docker, Linux, Windows o MacOS)
mediante el archivo de inicialización \texttt{init.sql} (o \texttt{seed.sql})
disponible en el repositorio oficial de \texttt{Apimetro} en GitHub.
